\documentclass[12pt]{article}
\usepackage[utf8]{inputenc}
\usepackage[T1]{fontenc}
\usepackage[margin=1in]{geometry}
\usepackage{setspace}
\usepackage{xcolor}
\usepackage[colorlinks=true, urlcolor=blue!60!black, linkcolor=blue!60!black]{hyperref}
\setlength{\parindent}{2em}
\setlength{\parskip}{1.25em}
\renewcommand{\baselinestretch}{1.0}
\begin{document}
\begin{center}
{\LARGE \textit{Whole Woman's Health v. Jackson}, 595 U.S. 30 (2021)\par}
\vspace{0.8em}
{\large Case Brief\par}
\end{center}
\vspace{1em}

\section{Memory Jogger}

State sovereign immunity and traditional equitable limits under \textit{Ex parte Young} barred federal pre-enforcement suits against state judges, court clerks, and the Texas Attorney General to block Texas's private-enforcement abortion statute (S.B. 8), leaving only a narrow path against state executive licensing officials (\href{https://www.courtlistener.com/opinion/5305640/whole-womans-health-v-jackson/}{CourtListener}; \href{https://www.supremecourt.gov/opinions/21pdf/21-463_3ebh.pdf}{Slip Opinion}).

\section{Detailed Case Facts}

In 2021, Texas enacted Senate Bill 8 (``S.B. 8''), titled the Texas Heartbeat Act and codified at Tex. Health \& Safety Code Ann. \S\S~171.201--171.212, effective September 1, 2021 (\href{https://supreme.justia.com/cases/federal/us/595/21-463/}{Justia}). The statute prohibited physicians from knowingly performing or inducing an abortion once a fetal heartbeat was detected---roughly six weeks after a patient's last menstrual period, well prior to fetal viability---unless a medical emergency prevented compliance. At the time, that restriction directly conflicted with \textit{Roe v. Wade}, 410 U.S. 113 (1973), and \textit{Planned Parenthood of Southeastern Pennsylvania v. Casey}, 505 U.S. 833 (1992).

To thwart ordinary federal pre-enforcement review, the Texas Legislature expressly forbade state and local officials from enforcing S.B. 8 directly through criminal prosecutions or civil actions (\S~171.207(a)). Instead, the statute delegated enforcement exclusively to private citizens, authorizing any non-governmental person to sue anyone who performed, induced, aided or abetted, or intended to aid or abet a prohibited abortion (\S~171.208(a)). Prevailing private plaintiffs could recover injunctive relief, statutory damages of at least \$10,000 per violation, plus costs and attorney's fees (\S~171.208(b)).

S.B. 8 also stacked state procedural rules heavily against defendants: plaintiffs could sue in their own home county among Texas's 254 counties and veto venue transfers (\S~171.210); prevailing defendants could not recover attorney's fees (\S~171.208(i)); nonmutual issue and claim preclusion were barred, exposing defendants to serial suits over a single procedure (\S~171.208(e)(5)); the statutory ``undue burden'' defense was tightly cabined (\S~171.209); and liability could attach retroactively if a preliminary injunction or \textit{Roe}/\textit{Casey} were later reversed (\S\S~171.208(e)(3), 171.209(e)) (\href{https://supreme.justia.com/cases/federal/us/595/21-463/}{Justia}).

Before S.B. 8 took effect, a group of Texas abortion providers and clinics led by Whole Woman's Health brought a federal pre-enforcement action under 42 U.S.C. \S~1983 seeking declaratory and injunctive relief. They named four categories of defendants: (1) Texas state district judge Austin Reeve Jackson and state-court clerk Penny Clarkston (as putative class representatives of all Texas state judges and clerks); (2) Texas Attorney General Ken Paxton; (3) four state executive licensing officials---Stephen Carlton (Texas Medical Board), Katherine Thomas (Texas Board of Nursing), Allison Benz (Texas Board of Pharmacy), and Cecile Young (Texas Health and Human Services Commission); and (4) Mark Lee Dickson, a private anti-abortion activist (\href{https://www.supremecourt.gov/opinions/21pdf/21-463_3ebh.pdf}{Slip Opinion}).

\section{Procedural History}

\begin{itemize}
    \item \textbf{U.S. District Court (W.D. Tex.):} On August 25, 2021, Judge Robert Pitman denied the state officials' Rule 12(b)(1) motions to dismiss on sovereign immunity and Article III grounds, as well as Mr. Dickson's standing challenge (\href{https://www.supremecourt.gov/opinions/21pdf/21-463_3ebh.pdf}{Slip Opinion}).
    \item \textbf{U.S. Court of Appeals (5th Cir.):} The state officials took an immediate interlocutory appeal under the collateral order doctrine, and the Fifth Circuit also entertained Mr. Dickson's appeal. The Fifth Circuit refused to enjoin S.B. 8 pending appeal and stayed all district court proceedings (\href{https://supreme.justia.com/cases/federal/us/595/21-463/}{Justia}).
    \item \textbf{Emergency Application (Sept. 1, 2021):} On the night S.B. 8 took effect, the Supreme Court denied emergency injunctive relief by a 5--4 vote in an unsigned shadow-docket order, citing unresolved procedural and sovereign-immunity questions (\href{https://www.law.cornell.edu/supremecourt/text/21A24}{Cornell LII, No. 21A24}).
    \item \textbf{Certiorari Before Judgment:} On October 22, 2021, the Supreme Court granted petitioners' expedited request for certiorari before judgment, heard oral argument on November 1, 2021, and issued its decision on December 10, 2021 (\href{https://www.supremecourt.gov/opinions/21pdf/21-463_3ebh.pdf}{Slip Opinion}).
    \item \textbf{Subsequent History on Remand:} After the Supreme Court denied mandamus to force immediate remand to the district court (\href{https://www.law.cornell.edu/supremecourt/text/21-962}{\textit{In re Whole Woman's Health}, Jan. 20, 2022}), the Fifth Circuit certified the state-law enforcement question to the Supreme Court of Texas. On March 11, 2022, the Texas Supreme Court unanimously held that Texas law barred the licensing officials from enforcing S.B. 8 directly or indirectly, ending the federal suit (\href{https://www.txcourts.gov/media/1453763/220033.pdf}{Tex. Sup. Ct. Opinion}). Three months later, \href{https://www.supremecourt.gov/opinions/21pdf/19-1392_6j37.pdf}{\textit{Dobbs v. Jackson Women's Health Organization}} (2022) overruled \textit{Roe} and \textit{Casey}.
\end{itemize}

\section{Judicial Votes}

The Court fractured sharply across the four categories of defendants (\href{https://www.supremecourt.gov/opinions/21pdf/21-463_3ebh.pdf}{Slip Opinion}; \href{https://www.courtlistener.com/opinion/5305640/whole-womans-health-v-jackson/}{CourtListener}):

\begin{itemize}
    \item \textbf{Opinion of the Court (Parts I, II-A, II-B, II-D, III, IV):} Authored by Justice Gorsuch and joined by Justices Thomas, Alito, Kavanaugh, and Barrett (5--4 as to court clerks and Attorney General Paxton; 9--0 in judgment as to Judge Jackson and private defendant Dickson).
    \item \textbf{Plurality Opinion (Part II-C):} Authored by Justice Gorsuch and joined by Justices Alito, Kavanaugh, and Barrett. Combined with the four justices concurring in the judgment in part (Roberts, Breyer, Sotomayor, and Kagan), \textbf{8 justices} agreed that claims against the four executive licensing officials could survive a motion to dismiss.
    \item \textbf{Concurring in Part and Dissenting in Part:} Justice Thomas joined all of Justice Gorsuch's opinion except Part II-C. He was the lone dissenter (1--8) against allowing suit to proceed against the licensing officials.
    \item \textbf{Concurring in the Judgment in Part and Dissenting in Part:} Chief Justice Roberts, joined by Justices Breyer, Sotomayor, and Kagan (4 justices), agreed that claims against the licensing officials should proceed and that state judges (and Dickson) should be dismissed, but dissented (4--5) from dismissing state-court clerks and Attorney General Paxton.
    \item \textbf{Concurring in the Judgment in Part and Dissenting in Part:} Justice Sotomayor, joined by Justices Breyer and Kagan (3 justices), wrote separately to emphasize the systemic threat posed by S.B. 8's procedural design to constitutional supremacy and federal rights.
\end{itemize}

\section{Holding}

\textbf{Judgment:} Affirmed in part, reversed in part, and remanded. The Court directed dismissal of Judge Jackson (9--0), Clerk Clarkston (5--4), Attorney General Paxton (5--4), and private defendant Mark Lee Dickson (9--0), while permitting the pre-enforcement suit to proceed past the Rule 12(b)(1) stage against the four executive licensing officials (8--1) (\href{https://www.supremecourt.gov/opinions/21pdf/21-463_3ebh.pdf}{Slip Opinion}).

\textbf{Rule Announced:} Under the Eleventh Amendment, state sovereign immunity, and Article III, the \textit{Ex parte Young} exception to sovereign immunity is bounded by traditional equity practice. It permits pre-enforcement federal injunctive suits only against state executive officials who possess specific authority to enforce the challenged statute (or collateral disciplinary provisions triggered by it). It does not authorize federal injunctions restraining state-court judges from adjudicating private civil suits, state-court clerks from docketing filings, or a state attorney general who lacks statutory enforcement authority over the challenged law; nor may a federal court enjoin unnamed private citizens or ``laws themselves'' (\href{https://supreme.justia.com/cases/federal/us/595/21-463/}{Justia}).

\textbf{Holding vs. Dicta:}
\begin{itemize}
    \item \textbf{Binding Holdings:} (1) Sovereign immunity and Article III's adversity requirement bar pre-enforcement federal suits seeking to enjoin state-court judges and clerks from adjudicating or docketing private civil actions under S.B. 8; (2) an attorney general lacking statutory enforcement power over the challenged subtitle cannot be sued under \textit{Ex parte Young}; (3) at the pleading stage, executive licensing officials with apparent statutory authority to discipline licensees for Chapter 171 violations fall within \textit{Ex parte Young}; and (4) a private individual who submits uncontested sworn declarations disclaiming any intent to sue cannot be sued for lack of Article III standing.
    \item \textbf{Dicta:} Justice Gorsuch's Part III discussion assuring litigants that state-court pre-enforcement suits, defensive litigation, and potential congressional legislation provide adequate alternative avenues to vindicate federal supremacy was unnecessary to the jurisdictional disposition. Likewise, his tentative reading of Texas's saving clause in Part II-C explicitly acknowledged that Texas courts remained the final arbiters of state statutory meaning---a reservation the Texas Supreme Court later invoked to extinguish the remaining claims (\href{https://www.txcourts.gov/media/1453763/220033.pdf}{Tex. Sup. Ct. Opinion}).
\end{itemize}

\section{Analysis of Opinions}

\subsection{Majority and Plurality Opinion (Justice Gorsuch)}

Writing for a five-justice majority (except as to Part II-C), Justice Gorsuch framed the appeal strictly around federal jurisdiction and sovereign immunity rather than S.B. 8's substantive constitutionality. Grounding state immunity in the Eleventh Amendment and broader constitutional structure (\href{https://supreme.justia.com/cases/federal/us/527/706/}{\textit{Alden v. Maine}}, 527 U.S. 706 (1999)), he emphasized that \textit{Ex parte Young}, 209 U.S. 123 (1908), carved out only a narrow equitable exception allowing injunctions against state executive officers who enforce unconstitutional laws.

Applying that framework in Part II-A, Justice Gorsuch rejected petitioners' attempt to enjoin state judges and clerks. Unlike executive enforcers, judges and clerks resolve or process disputes filed by private litigants; \textit{Ex parte Young} itself warned that enjoining a state court or its machinery violates the constitutional scheme of dual sovereignty. Furthermore, Article III requires an actual controversy between adverse litigants, and clerks who perform ministerial docketing duties under Texas Rule of Civil Procedure 24 are neither adverse to prospective defendants nor authorized to screen filings on the merits. Permitting clerk-targeted injunctions would lack any workable limiting principle and invite federal supervision of state judicial operations (\href{https://supreme.justia.com/cases/federal/us/595/21-463/}{Slip Opinion}).

In Part II-B, the majority dismissed Attorney General Paxton because petitioners failed to identify any statutory enforcement authority he held over S.B. 8. Texas Occupations Code \S~165.101 authorizes civil-penalty suits only for violations of that subtitle or a Texas Medical Board rule, whereas S.B. 8 resides in the Health and Safety Code. Even if Paxton held some enforcement power, traditional equity forbids courts from enjoining ``the world at large'' or binding unnamed private litigants through an injunction against a single state official.

In Part II-C (for a four-justice plurality), Justice Gorsuch concluded that the four executive licensing officials fell within \textit{Ex parte Young} at the motion-to-dismiss stage. Although S.B. 8 purported to make private suits exclusive, its saving clause (\S~171.207(b)(3)) preserved enforcement of other laws regulating abortion, and Texas Occupations Code \S~164.055 directed the Texas Medical Board to discipline physicians violating Chapter 171 of the Health and Safety Code. In Part II-D, the Court unanimously dismissed Mr. Dickson because his uncontested sworn disclaimer negated Article III traceability. Finally, in Part III, the majority pushed back against Justice Sotomayor's warnings of nullification, noting that the Constitution guarantees no universal right to federal pre-enforcement review, that federal defenses remain fully available in state-court actions under the Supremacy Clause, and that ``chilling effect'' alone cannot override Article III and equitable limits.

\subsection{Concurring in Part and Dissenting in Part (Justice Thomas)}

Justice Thomas joined every part of Justice Gorsuch's opinion except Part II-C and would have ordered dismissal of all defendants (\href{https://supreme.justia.com/cases/federal/us/595/21-463/}{Justia}). Relying on historical equity practice exercised by the English High Court of Chancery, he contended that \textit{Ex parte Young} authorizes only a ``negative injunction'' against an officer who both possesses statutory enforcement power and poses an imminent threat of commencing proceedings. Reading S.B. 8 as a harmonious whole, Justice Thomas argued that \S~171.207(a) ``triply'' stripped all state officials---including licensing boards---of enforcement authority, and that the saving clause preserved only independent substantive abortion bans, not collateral licensing enforcement of S.B. 8. Moreover, petitioners alleged only a speculative chill rather than a credible, imminent threat of license revocation. Notably, the Texas Supreme Court later adopted Justice Thomas's reading of Texas law (\href{https://www.txcourts.gov/media/1453763/220033.pdf}{Tex. Sup. Ct. Opinion}).

\subsection{Concurring in the Judgment in Part and Dissenting in Part (Chief Justice Roberts)}

Chief Justice Roberts, joined by Justices Breyer, Sotomayor, and Kagan, agreed that the licensing officials were proper defendants and that state judges were not, because judges are not adverse to litigants (\href{https://www.supremecourt.gov/opinions/21pdf/21-463_3ebh.pdf}{Slip Opinion}). However, he argued that both Attorney General Paxton and Clerk Clarkston should also remain in the suit. Under Texas Administrative Code \S~190.8(7) and Occupations Code \S~165.101, the Attorney General held coextensive civil-penalty authority to enforce Medical Board rules incorporating Chapter 171. As for court clerks, Chief Justice Roberts reasoned that S.B. 8's punitive procedural design meant the mere filing and docketing of suits inflicted the chilling injury, making clerks who set the court's machinery in motion sufficiently connected to enforcement under \textit{Ex parte Young}. Invoking \textit{Marbury v. Madison} and \textit{United States v. Peters}, he warned that allowing a state to nullify Supreme Court precedent through procedural artifice threatened the constitutional role of the Supreme Court itself, regardless of the underlying right at stake.

\subsection{Concurring in the Judgment in Part and Dissenting in Part (Justice Sotomayor)}

Justice Sotomayor, joined by Justices Breyer and Kagan, delivered a forceful dissent emphasizing the practical devastation S.B. 8 had already inflicted on abortion access in Texas (\href{https://supreme.justia.com/cases/federal/us/595/21-463/}{Justia}). She argued that \textit{Ex parte Young} was created precisely to prevent states from using ruinous penalties to coerce compliance and deny any meaningful judicial hearing. Unlike ordinary tort laws, S.B. 8 deputized uninjured strangers as bounty hunters while dismantling venue fairness, fee-shifting parity, and preclusion rules to turn state courts into a harassment weapon. Citing \href{https://supreme.justia.com/cases/federal/us/563/247/}{\textit{Virginia Office for Protection and Advocacy v. Stewart}}, 563 U.S. 247 (2011), she maintained that \textit{Ex parte Young} should adapt to novel evasions rather than reward them. Comparing Texas's strategy to John C. Calhoun's nullification doctrine, she warned that the majority's roadmap invited other states to copy and perfect S.B. 8 to chill Second Amendment, religious, or free-speech rights simply by scrubbing all residual executive licensing hooks.

\section{Examples: Future Applications}

\begin{enumerate}
    \item \textbf{Same-Side (Barred by Sovereign Immunity):} A state enacts a statute banning private possession of certain semiautomatic firearms, expressly prohibits all state and local officials from enforcing it directly or through licensing boards, and authorizes any private citizen to sue gun owners for \$15,000 in statutory damages. Gun owners sue the state's chief trial judge and county court clerks in federal court to block docketing of private suits. Under \textit{Whole Woman's Health}, the suit must be dismissed because state judges and clerks are immune under the Eleventh Amendment and lack Article III adversity.
    \item \textbf{Same-Side (Barred by Sovereign Immunity):} A state passes a private-bounty statute targeting social media platforms that moderate political speech, explicitly stripping the state attorney general of enforcement authority. A platform sues the attorney general in federal court, arguing that an injunction against him will bind private citizens statewide. Under \textit{Whole Woman's Health}, the attorney general must be dismissed for lack of statutory enforcement connection, and federal courts cannot enjoin unnamed private litigants.
    \item \textbf{Opposite-Side (Permitted Under \textit{Ex parte Young}):} A state enacts a private-enforcement statute penalizing religious schools that teach disfavored doctrines, but a separate provision of the state education code requires the State Commissioner of Education to revoke the accreditation of any school that violates the statute. A religious school sues the Commissioner in federal court to enjoin accreditation revocation. Under Part II-C of \textit{Whole Woman's Health}, the suit may proceed past a motion to dismiss because an executive official retains specific collateral enforcement authority.
    \item \textbf{Opposite-Side (Permitted Against Private Enforcer):} Under an S.B. 8-style private-bounty law restricting drag performances, a named private activist sends certified demand letters to a specific theater owner threatening to file a \$10,000 bounty suit on Monday unless the weekend show is canceled. The theater owner sues that named activist in federal court for declaratory and injunctive relief. Unlike Mr. Dickson---who disclaimed any intent to sue---this private defendant's concrete, imminent threat establishes Article III injury, traceability, and redressability, and private citizens enjoy no Eleventh Amendment immunity.
    \item \textbf{Fence-Sitter (Close Case):} A state enacts a private-bounty statute banning gender-affirming medical care and expressly forbids executive enforcement, but requires a state Department of Health registrar to review, stamp, and publish a mandatory ``Certificate of Compliance'' before any private plaintiff may file a bounty suit in state court. Plaintiffs sue the registrar in federal court. On one hand, the registrar is an executive officer whose affirmative gatekeeping act is a statutory prerequisite to private enforcement, distinguishing them from passive court clerks; on the other hand, defendants would argue the registrar's role is purely ministerial and analogous to docketing under \textit{Whole Woman's Health}, leaving federal courts closely divided on whether \textit{Ex parte Young} applies.
\end{enumerate}

\section{Critique}

\textit{Whole Woman's Health v. Jackson} sparked intense debate across federal courts scholarship. From an \textbf{originalist and formalist perspective}, defenders praised the majority's fidelity to traditional equitable boundaries, Article III's case-or-controversy requirement, and the structural dignity of state judiciaries. Because federal equity jurisdiction tracks historic Chancery practice, federal courts lack roving authority to invent novel structural remedies against state court personnel or to issue injunctions binding the general public. Indeed, Justice Thomas's textualist critique of Part II-C proved prescient when the Texas Supreme Court confirmed that Texas law barred licensing enforcement as well (\href{https://texaslawreview.org/judicial-review-of-unconventional-enforcement-regimes/}{Texas Law Review}).

From a \textbf{progressive and doctrinal-remedies perspective}, scholars sharply criticized the decision as an abdication of the federal judiciary's duty to enforce constitutional supremacy. Richard Fallon argued in the \textit{Harvard Law Review} that the Court's rigid separation-of-powers formalism wrongly truncated constitutional remedies, cheapening substantive rights whenever legislatures devise novel evasion tactics (\href{https://harvardlawreview.org/print/vol-136/constitutional-remedies-in-one-era-and-out-the-other/}{Fallon, 136 Harv. L. Rev. 1300 (2023)}). Similarly, Jon Michaels and David Noll warned that the ruling catalyzed ``vigilante federalism,'' encouraging states across the political spectrum---including California's S.B. 8-inspired firearm legislation---to outsource enforcement of constitutionally suspect laws to private bounty hunters (\href{https://publications.lawschool.cornell.edu/lawreview/wp-content/uploads/sites/2/2023/09/Michaels-Noll-final.pdf}{Michaels \& Noll, 108 Cornell L. Rev. 1187 (2023)}). Subsequent commentary also observed that \textit{Whole Woman's Health} created doctrinal tension with \textit{Pulliam v. Allen} by erecting a formidable new sovereign-immunity barrier to suits against state judges (\href{https://harvardlawreview.org/wp-content/uploads/2026/05/139-Harv.-L.-Rev.-1697.pdf}{139 Harv. L. Rev. 1697}).

At bottom, the case reflects a deep clash of \textbf{competing constitutional values}: state sovereign immunity, judicial restraint, and traditional limits on federal equitable power on one side, versus the Supremacy Clause, the availability of effective remedies for constitutional violations, and the prevention of legislative gamesmanship designed to chill federal rights on the other.

\section{Key Quotations}

\begin{quote}
``As \textit{Ex parte Young} put it, `an injunction against a state court' or its `machinery' `would be a violation of the whole scheme of our Government.''' \\
--- Justice Gorsuch, Opinion of the Court (\href{https://www.supremecourt.gov/opinions/21pdf/21-463_3ebh.pdf}{Slip Opinion}; \href{https://www.courtlistener.com/opinion/5305640/whole-womans-health-v-jackson/}{CourtListener})
\end{quote}

\begin{quote}
``Consistent with historical practice, a federal court exercising its equitable authority may enjoin named defendants from taking specified unlawful actions. But under traditional equitable principles, no court may `lawfully enjoin the world at large,' or purport to enjoin challenged `laws themselves.''' \\
--- Justice Gorsuch, Opinion of the Court (\href{https://supreme.justia.com/cases/federal/us/595/21-463/}{Justia}; \href{https://www.courtlistener.com/opinion/5305640/whole-womans-health-v-jackson/}{CourtListener})
\end{quote}

\begin{quote}
``The clear purpose and actual effect of S. B. 8 has been to nullify this Court's rulings. . . . The nature of the federal right infringed does not matter; it is the role of the Supreme Court in our constitutional system that is at stake.'' \\
--- Chief Justice Roberts, concurring in the judgment in part and dissenting in part (\href{https://www.supremecourt.gov/opinions/21pdf/21-463_3ebh.pdf}{Slip Opinion})
\end{quote}

\begin{quote}
``By foreclosing suit against state-court officials and the state attorney general, the Court effectively invites other States to refine S. B. 8's model for nullifying federal rights.'' \\
--- Justice Sotomayor, concurring in the judgment in part and dissenting in part (\href{https://supreme.justia.com/cases/federal/us/595/21-463/}{Justia})
\end{quote}

\end{document}
